\input{estilo}
% El pipeline desde la fuente, con el watermark de cada etapa en rojo y las
% etapas a la misma distancia. Cada etapa atrasa el watermark lo que tarda en
% procesar lo que tiene pendiente (los 10 minutos del ejemplo): por eso
% W(A) ≥ W(B) ≥ W(C).
\begin{document}
\begin{figura}\begin{tikzpicture}
  \node[caja, draw=suave, fill=white, text=suave] (f) at (0,0) {fuente};
  \foreach \n/\x/\h in {a/2.8/8:30, b/5.6/8:20, c/8.4/8:10} {
    \node[caja] (\n) at (\x,0) {\MakeUppercase{\n}};
    \node[etiq, text=acento, above=2pt of \n] {\var{W} = \h};
  }
  \draw[msg] (f) -- (a);
  \draw[msg] (a) -- node[etiq, text=suave, below=3pt] {−10 min} (b);
  \draw[msg] (b) -- node[etiq, text=suave, below=3pt] {−10 min} (c);
  \node[anchor=north west, align=left, font=\small] at (-0.55,-1.2) {%
    \var{W}(A) ≥ \var{W}(B) ≥ \var{W}(C)\\[4pt]
    \var{W}(B) = min (\var{W}(A), trabajo más viejo pendiente en B)};
\end{tikzpicture}\end{figura}
\end{document}
